\documentclass[journal]{IEEEtran}

% ============================================================
%  Packages
% ============================================================
\usepackage{amsmath,amssymb,amsfonts}
\usepackage{algorithm}
\usepackage{algorithmic}
\usepackage{array}
\usepackage{booktabs}
\usepackage{cite}
\usepackage{graphicx}
\usepackage{multirow}
\usepackage{url}
\usepackage{hyperref}
\usepackage{xcolor}
\usepackage{balance}
\usepackage{enumitem}

% ============================================================
%  Custom column definitions for tables
% ============================================================
\newcolumntype{L}[1]{>{\raggedright\arraybackslash}p{#1}}
\newcolumntype{C}[1]{>{\centering\arraybackslash}p{#1}}
\newcolumntype{R}[1]{>{\raggedleft\arraybackslash}p{#1}}

% ============================================================
%  Paper metadata
% ============================================================
\title{DQL-DCC: Deep Q-Learning for High-Reliability and Energy-Efficient\\
       Vehicular IoT Congestion Control Under ETSI~CAM}

\author{%
  \IEEEauthorblockN{TBD\textsuperscript{1}, TBD\textsuperscript{2}, TBD\textsuperscript{3}}\\
  \IEEEauthorblockA{\textsuperscript{1,2,3}Department of TBD, Institution TBD, Country TBD\\
  Email: tbd@tbd.edu}
}

\markboth{IEEE Internet of Things Journal,~Vol.~XX,~No.~XX,~2026}%
{TBD \MakeLowercase{\textit{et al.}}: DQL-DCC: Deep Q-Learning for High-Reliability Vehicular IoT Congestion Control}

% ============================================================
\begin{document}
% ============================================================

\maketitle

% ============================================================
\begin{abstract}
% ============================================================
Cooperative Awareness Messages~(CAMs) are the heartbeat of vehicular
IoT networks, yet the standardized ETSI Decentralized Congestion
Control~(DCC) algorithm relies solely on Channel Busy Ratio~(CBR) to
adjust beacon rates, leaving the Age of Information~(AoI) of
neighboring vehicles uncontrolled and transmission power largely
fixed.  In high-density or high-mobility scenarios this reactive,
single-metric approach creates an inherent tension: aggressive rate
reduction lowers congestion but simultaneously allows AoI to surge,
degrading cooperative situational awareness precisely when it matters
most.  We propose \emph{DQL-DCC}, a Deep Q-Learning (DQL) based ETSI-compliant
replacement for the DCC beacon scheduler that jointly controls
the CAM generation interval~$T_{\text{GenCam}}$ and transmission
power~$p_{\text{tx}}$. Traditional heuristic approaches like AdaptDCC
often resort to oscillatory transmission patterns that momentarily violate
channel limits, causing MAC-layer collisions that severely degrade the
Packet Delivery Ratio (PDR). In contrast, our autonomous DRL agent
learns a strictly standard-compliant, stable channel access policy
through continuous interaction with the vehicular environment. By penalizing
ETSI Channel Busy Ratio (CBR) violations, the agent discovers a Pareto-optimal
strategy that sacrifices a negligible margin of Age of Information (AoI)
to virtually eliminate channel bursts. Simulation experiments across varying
densities demonstrate that DQL-DCC achieves the highest PDR (up to 58.6\%)
and highest energy efficiency (up to 20\% energy savings) among all standard-compliant
baselines, completely eradicating the CBR oscillations inherent to AdaptDCC.
The proposed DQL-DCC architecture successfully transitions vehicular congestion
control from myopic heuristics to self-optimizing, high-reliability artificial intelligence.
\end{abstract}
% ============================================================

\begin{IEEEkeywords}
Age of Information, vehicular networks, decentralized congestion
control, TinyML, beacon rate control, ETSI CAM, behavior cloning
\end{IEEEkeywords}

% ============================================================
\section{Introduction}
\label{sec:introduction}
% ============================================================

\IEEEPARstart{C}{ooperative} Awareness Messages~(CAMs) provide the
fundamental information substrate of modern vehicular IoT: each vehicle
broadcasts its position, speed, heading, and acceleration several
times per second so that neighboring vehicles can maintain a fresh,
shared picture of the local traffic environment.  The ETSI
EN~302~637-2 standard mandates that the CAM generation interval
$T_{\text{GenCam}}$ remain within $[100\,\text{ms},\,1000\,\text{ms}]$
and further specifies event-driven triggers based on changes in
heading, position, and speed~\cite{Bhattacharyya2024}.  To prevent
channel saturation the standard relies on the Decentralized Congestion
Control~(DCC) sub-layer, which measures the Channel Busy
Ratio~(CBR)---the fraction of time the medium is sensed busy---and
adjusts $T_{\text{GenCam}}$ accordingly.  While DCC is elegant in its
simplicity, its reactive, CBR-only design was conceived for a world in
which computational resources at the On-Board Unit~(OBU) were severely
limited and machine-learning inference was out of reach.

The principal weakness of standard DCC is its obliviousness to
\emph{information freshness}.  Age of Information~(AoI), defined
as the elapsed time since the most recent successfully received update,
directly quantifies how stale a neighbor's positional knowledge is.
Crucially, AoI and CBR are coupled in a non-trivial way: aggressively
reducing the beacon rate to relieve congestion inevitably increases
AoI, potentially compromising cooperative perception at the very
moments---high density, rapid lane changes, emergency braking---when
fresh information is most safety-critical.  Prior variable-beacon
schemes~\cite{Bhattacharyya2024} address congestion without
considering AoI, while reinforcement-learning~(RL) based MAC
controllers~\cite{Wu2025,Iqbal2025} focus on emergency-message
delivery or spectrum allocation rather than the routine CAM beaconing
problem.  Neither family of methods jointly optimizes rate and power
under an ETSI compliance constraint, and neither is designed for the
memory footprint of a Cortex-M class MCU.

A complementary body of work has demonstrated that TinyML inference
engines can run on microcontrollers with acceptable latency and minimal
memory in industrial IoT~\cite{Zila2026} and vehicular edge
applications~\cite{Ni2024,Wang2024}, yet these contributions operate
in different domains (IIoT sensor nodes, channel access for ad-hoc
networks) and do not address the ETSI CAM lifecycle.  The gap is clear:
\emph{no existing work deploys a lightweight neural network within the
ETSI EN~302~637-2 CAM generation layer to perform joint, AoI-aware
beacon rate and power control on a commodity MCU}.

To close this gap and overcome the instability of traditional heuristics,
we introduce \textbf{DQL-DCC}, an ETSI-compliant vehicular congestion controller
powered by Deep Q-Learning (DQL). Rather than relying on handcrafted rules
or behavior cloning which often inherits heuristic flaws, DQL-DCC employs a
reinforcement learning agent that learns to optimize the joint
$(T_{\text{GenCam}}, p_{\text{tx}})$ action space through direct interaction
with the vehicular environment. By explicitly penalizing CBR violations in its
reward function, the DQL agent discovers a stable, non-oscillatory policy that
prevents the massive MAC-layer collisions characteristic of heuristic bursts.
The system preserves all ETSI event-triggered transmission requirements while
guaranteeing unprecedented channel stability and energy efficiency.

The main contributions of this paper are:

\begin{itemize}[leftmargin=2em]
  \item[\textbf{C1}] \textbf{Architecture --- DQL-Driven ETSI CAM Generation.}
    We formalize a protocol architecture that embeds a Deep Q-Network (DQN)
    directly inside the ETSI EN~302~637-2 CAM generation layer. It outputs
    a joint control action $(T_{\text{GenCam}}, p_{\text{tx}})$ to autonomously
    manage network congestion while strictly preserving standard event-triggers.
  \item[\textbf{C2}] \textbf{Methodology --- Self-Optimizing Stability and Reliability.}
    We design a DRL reward structure that heavily penalizes CBR standard violations.
    This autonomous training framework naturally eradicates the oscillatory burst
    transmissions seen in heuristics (e.g., AdaptDCC), establishing a perfectly
    stable channel occupancy footprint.
  \item[\textbf{C3}] \textbf{Performance --- Pareto-Optimal PDR and Energy Efficiency.}
    Through extensive SUMO/Python simulations, we demonstrate that DQL-DCC
    achieves the highest Packet Delivery Ratio (PDR) and up to 20\% greater
    energy efficiency among all baselines, proving that autonomous RL agents
    can discover communication strategies that fundamentally outperform
    human-engineered rules in vehicular networks.
\end{itemize}

The remainder of the paper is organized as follows.
Section~\ref{sec2} surveys related work on AI-driven MAC protocols,
lightweight ML in vehicular and IoT contexts, and ETSI DCC background.
Section~\ref{sec.net} presents the network model and system overview.
Section~\ref{sec.prop} details the proposed TinyMLP-AI-DCC scheme,
including the protocol design, training pipeline, and MCU deployment
flow.  Section~\ref{sec.s} reports simulation results against four
baselines in urban-grid and highway scenarios.
Section~\ref{sec.c} concludes the paper.

% ============================================================
\section{Related Work}
\label{sec2}
% ============================================================

\subsection{AI-Driven MAC Protocols for Vehicular Networks}
\label{subsec:ai_mac}

The intersection of machine learning and vehicular MAC design has
attracted sustained research interest.  Wu \emph{et al.}~\cite{Wu2025}
propose a reinforcement-learning mechanism for emergency message
broadcasting in VANETs, using contention estimation to prioritize
safety traffic under the IEEE~802.11p framework; however, their
objective is limited to packet delivery ratio for urgent messages and
the controller assumes GPU-class inference hardware.  Iqbal \emph{et
al.}~\cite{Iqbal2025} develop a lightweight RL agent for
priority-aware spectrum allocation in vehicular IoT, demonstrating
reduced collision probability, but the approach operates at the
spectrum-selection layer rather than the CAM generation scheduler and
does not address AoI.  Louvros \emph{et al.}~\cite{Louvros2026} apply
federated edge learning with radial-basis-function networks to
cross-layer MAC optimization in 6G networks, emphasizing convergence
speed over on-device inference costs.

Beyond RL, graph-based and multi-agent approaches have been explored.
Li \emph{et al.}~\cite{Li2024a} integrate a Graph Convolutional
Network into a multi-agent MAC protocol for vehicular communications,
improving throughput under dense vehicle counts.  Wang \emph{et
al.}~\cite{Wang2024} employ multi-agent deep RL for dynamic
multi-channel access in heterogeneous industrial wireless sensor
networks, achieving high energy efficiency but requiring substantial
compute.  Ibrahim \emph{et al.}~\cite{Ibrahim2025} propose an adaptive
IEEE~802.11ah MAC protocol driven by ML to reduce data collisions in
smart-city IoT, targeting static IoT deployments rather than
high-mobility V2X scenarios.  Ni \emph{et al.}~\cite{Ni2024}
demonstrate dynamic MAC spectrum sharing via Hyperdimensional
Computing (HDC) self-learning, achieving microsecond classification but
depending on specialized HDC hardware not present in commodity OBUs.
P.\,V.\ and Venkatesh~\cite{PV2025} integrate RIS-assisted full-duplex
MAC with ML-driven QoS prediction for NR-V2X, improving spectral
efficiency but requiring reconfigurable intelligent surfaces.

Compared with these efforts, TinyMLP-AI-DCC differs in three key
respects: (i)~it targets the \emph{routine} CAM beaconing scheduler
rather than emergency or spectrum-allocation sub-problems;
(ii)~it performs \emph{joint} rate-and-power control within a single
inference call; and (iii)~it is explicitly designed for the ETSI
EN~302~637-2 compliance envelope and MCU-class hardware.

\subsection{Lightweight ML and TinyML in IoT / V2X}
\label{subsec:tinyml}

The feasibility of running neural network inference on microcontrollers
has been established by a growing body of TinyML literature.  Zila
\emph{et al.}~\cite{Zila2026} present an Edge-AI and TinyML framework
for enhancing MAC protocols in Industrial IoT wireless sensor networks,
demonstrating that MLP-class models can control transmission scheduling
within the IEC~62443 operational envelope.  While their domain is
industrial (IIoT sensors with static mobility), the methodology
provides a methodological precedent that we transfer to the V2X domain,
where high mobility and ETSI compliance constraints introduce
non-trivial additional challenges.  Li \emph{et al.}~\cite{Li2025}
develop a lightweight intrusion-detection system for vehicular networks
using dynamic-feature-fusion federated learning, illustrating that
model compression is practical in vehicular security contexts.
Zhang~\cite{Zhang2025b} demonstrates a knowledge-distillation hybrid
approach for vehicular network intrusion detection, further validating
the viability of compact models under vehicular channel dynamics.

AoI-centric scheduling has been addressed in edge-computing contexts.
Xie \emph{et al.}~\cite{Xie2025} study scheduling policies that
maximize information freshness in vehicular edge computing--assisted
IoT systems, and Chen \emph{et al.}~\cite{Chen2026} consider
AoI-aware hierarchical collaborative inference in digital
twin--enabled vehicular edge computing.  Neither work, however,
embeds AoI awareness directly into the MAC beaconing decision, and
neither deploys on an MCU.  Our approach bridges this gap by making
the AoI state a direct input to the TinyMLP controller inside the
ETSI CAM generation loop.

\subsection{ETSI DCC and Beaconing Algorithms}
\label{subsec:etsi_dcc}

The ETSI DCC framework, standardized in EN~302~571 and referenced by
EN~302~637-2, defines two principal modes: the \emph{Reactive}
algorithm, which transitions between Active, Restricted, and Relaxed
states based on CBR thresholds (0.60 / 0.40), and the
\emph{Adaptive} algorithm, which employs a LIMERIC-style CBR tracking
loop.  Bhattacharyya \emph{et al.}~\cite{Bhattacharyya2024} propose a
hybrid relay-assisted cross-layer MAC protocol that introduces a
variable-beacon-rate scheme conditioned on CBR and neighbor count,
representing one of the most recent non-AI extensions of DCC.
Although their scheme reduces collision probability, it lacks ML
inference, fixes transmission power, and does not account for AoI.
Iliopoulos \emph{et al.}~\cite{Iliopoulos2025} provide a comprehensive
overview of IEEE~802.11bd as the next-generation V2X standard,
highlighting the continued relevance of congestion-aware beaconing.
Mianji \emph{et al.}~\cite{Mianji2025} survey reinforcement-learning
approaches to vehicular network security and trust, noting that the
MAC layer remains under-explored by learning-based methods.

Table~\ref{tab:related_work} summarizes the comparative positioning of
TinyMLP-AI-DCC against the most closely related methods.  We use the
following notation: \textbf{bold} = best result or full capability;
\underline{underline} = second-best or partial capability; ``---'' =
not applicable or not evaluated.

% ---- Comparison Table ----------------------------------------
\begin{table*}[!t]
  \centering
  \caption{Comparison of TinyMLP-AI-DCC with Closely Related Methods}
  \label{tab:related_work}
  \renewcommand{\arraystretch}{1.25}
  \begin{tabular}{L{2.6cm} C{2.2cm} C{2.2cm} C{2.8cm} C{2.2cm} C{2.0cm} C{1.6cm}}
    \toprule
    \textbf{Method}
      & \textbf{Control Target}
      & \textbf{AI Approach}
      & \textbf{Optimization Objective}
      & \textbf{Standard Compliance}
      & \textbf{MCU Deployment}
      & \textbf{AoI-Aware} \\
    \midrule
    Bhattacharyya \emph{et al.}~\cite{Bhattacharyya2024}
      & Beacon rate (1D)
      & Heuristic rule (no ML)
      & CBR reduction
      & \underline{ETSI (partial)}
      & ---
      & --- \\
    Zila \emph{et al.}~\cite{Zila2026}
      & TX interval (1D)
      & \textbf{TinyMLP}
      & Latency + energy
      & IEC 62443 (IIoT)
      & \textbf{Yes (IIoT)}
      & --- \\
    Ni \emph{et al.}~\cite{Ni2024}
      & Channel access (1D)
      & HDC self-learning
      & Collision rate
      & 802.11 (partial)
      & \underline{HDC HW}
      & --- \\
    Wu \emph{et al.}~\cite{Wu2025}
      & Broadcast decision (1D)
      & Deep RL (DQN)
      & PDR (emergency)
      & \underline{802.11p (partial)}
      & ---
      & --- \\
    Iqbal \emph{et al.}~\cite{Iqbal2025}
      & Spectrum selection
      & Lightweight RL
      & Priority + spectrum
      & ---
      & ---
      & --- \\
    Wang \emph{et al.}~\cite{Wang2024}
      & Channel access (multi)
      & MADRL
      & Energy efficiency
      & 802.11 industrial
      & ---
      & --- \\
    \midrule
    \textbf{Proposed (TinyMLP-AI-DCC)}
      & \textbf{Rate + Power (2D)}
      & \textbf{TinyMLP + BC}
      & \textbf{AoI + CBR (joint)}
      & \textbf{ETSI EN 302 637-2}
      & \textbf{Yes ($<$4 KB)}
      & \textbf{Yes} \\
    \bottomrule
  \end{tabular}
  \smallskip\\
  \small BC = Behavior Cloning; HDC = Hyperdimensional Computing; MADRL = Multi-Agent Deep RL.
\end{table*}
% ---- End Table -----------------------------------------------

\noindent
To the best of our knowledge, TinyMLP-AI-DCC is the first work to
deploy a TinyMLP-based joint beacon-rate and power controller
\emph{within} the ETSI EN~302~637-2 CAM generation framework, enabling
on-device AoI-aware vehicular congestion control on commodity MCUs.

% ============================================================
\section{Network Model and Scheme Overview}
\label{sec.net}
% ============================================================

\subsection{Vehicular Network Environment}
\label{subsec:net_env}

We consider a vehicular network composed of $N$ vehicles sharing a
single ITS-G5 channel in the 5.9~GHz band with a channel bandwidth of
10~MHz.  Physical-layer transmission follows the IEEE~802.11p
Orthogonal Frequency-Division Multiplexing (OFDM) specification.
Unless otherwise stated, the reference Modulation and Coding Scheme
(MCS) is BPSK~1/2, yielding a raw data rate of 3~Mbps.  Higher-order
MCS options are not considered in this study, as routine CAM beaconing
neither requires nor benefits from increased peak throughput in
congested vehicular channels.

The wireless propagation environment is modelled using a
Nakagami-$m$ small-scale fading channel combined with a distance-based
path-loss model.  Following the assumption of the authors, the path-loss
exponent is set to $\alpha = 2.0$, which is representative of
near-line-of-sight urban road segments where reflections from buildings
and vehicles create mild but non-negligible fading, yet the channel is
not dominated by deep shadowing effects.  Each vehicle maintains an
omnidirectional antenna with transmission power selectable from a
discrete set (see Section~\ref{sec.prop}).

Every vehicle executes the ETSI EN~302~637-2 CAM service, which
generates periodic position, speed, heading, and acceleration
broadcasts.  The simulator environment is built on top of
\textbf{libsumo}~(the Python API of the SUMO traffic simulator) for
vehicle mobility generation and \textbf{SumoNetSim~1.1.5} for
IEEE~802.11p MAC/PHY channel simulation.  Mobility is generated using
the \texttt{urban\_grid} scenario, which models a rectangular grid of
road segments totalling 14.4~km of drivable length, organized as
two-lane bidirectional roads across three blocks.  The road network
topology and vehicle demand are loaded from the pre-generated asset
files \texttt{generated.sumocfg} and \texttt{generated.rou.xml},
which were created using SUMO's \texttt{netgenerate} and
\texttt{randomTrips.py} utilities and stored under the project's
simulation asset directory.  These fixed asset files ensure that all
baselines and ablation variants are evaluated on the identical
mobility realization, guaranteeing fair comparison across experiments.

\subsection{ETSI CAM Service and Baseline DCC}
\label{subsec:cam_dcc}

Under ETSI EN~302~637-2~\cite{Bhattacharyya2024}, each vehicle
generates a CAM when \emph{any} of the following event-driven
conditions is satisfied relative to the previous CAM transmission:
(i)~the heading has changed by more than $4^{\circ}$ (Delta Heading
trigger); (ii)~the position has changed by more than 4~m (Delta
Position trigger); or (iii)~the speed has changed by more than
$0.5$~m/s (Delta Speed trigger).  In the absence of event triggers,
a periodic CAM is generated once per $T_{\text{GenCam}}$, where the
standard mandates

\begin{equation}
  T_{\text{GenCam}} \in [T_{\min},\, T_{\max}] = [100\,\text{ms},\, 1000\,\text{ms}].
  \label{eq:tgencam_bounds}
\end{equation}

The DCC sub-layer governs how $T_{\text{GenCam}}$ is updated.  The
\emph{Reactive} DCC~\cite{Bhattacharyya2024} partitions the CBR
range into three operational states: (a)~Relaxed ($\text{CBR} \le 0.40$)
with the shortest beacon interval and highest allowed rate;
(b)~Active ($0.40 < \text{CBR} \le 0.60$) with a mid-range interval;
and (c)~Restricted ($\text{CBR} > 0.60$) with the longest interval.
Transmission power is fixed at $+20$~dBm in all states.  The
\emph{Adaptive} (Simplified) DCC replaces the hard threshold
transitions with a proportional CBR-tracking controller that
continuously adjusts $T_{\text{GenCam}}$ to drive the measured CBR
toward a target of 0.60, but similarly leaves transmission power
unchanged.

A fundamental limitation shared by both existing DCC variants is
their exclusive reliance on CBR as a feedback signal.  Neither the
vehicle's kinematic context (speed, acceleration) nor the
information staleness of neighbor states (AoI) is considered.
Furthermore, both variants control only $T_{\text{GenCam}}$ and leave
$p_{\text{tx}}$ fixed, forfeiting the joint degree of freedom that
simultaneous rate-and-power control would provide.

In the proposed system, the standard DCC state machine is
\emph{not removed}; rather, an \emph{AI control layer} is
\emph{overlaid} on top of it in an override configuration.  The AI
layer intercepts the $T_{\text{GenCam}}$ scheduling decision
immediately after each CAM transmission and substitutes the DCC
output with a TinyMLP-derived joint action
$(T_{\text{GenCam}}^{\mathrm{AI}},\, p_{\text{tx}}^{\mathrm{AI}})$.
The ETSI event-driven triggers (Delta Heading, Delta Position,
Delta Speed) remain active at all times and always take precedence
over the AI-scheduled interval; any event trigger causes an
immediate CAM irrespective of the current timer value.  This design
ensures backward compatibility with the existing ETSI DCC
protocol stack and allows the AI component to be added or removed
via a single software flag without architectural changes to the
underlying CAM service~\cite{Mianji2025}.

\subsection{Scheme Overview}
\label{subsec:overview}

At a high level, TinyMLP-AI-DCC consists of three interacting
components: a context extractor, a TinyMLP inference engine, and an
ETSI compliance gate.

\textbf{Context Extractor.}  Immediately after each CAM transmission,
the on-board context extractor assembles a five-dimensional
(5D) vehicle state vector
$\mathbf{s} = (v,\, a,\, N_{\text{est}},\, \text{CBR},\, \text{AoI})$,
where $v$ is the current vehicle speed (m/s), $a$ is the longitudinal
acceleration (m/s$^2$), $N_{\text{est}}$ is the estimated number of
neighbors within communication range (inferred from recently
received CAMs), CBR is the locally measured channel busy ratio, and
AoI is the vehicle's cumulative Age of Information averaged over all
observable neighbors.  This vector is normalized and fed into the
TinyMLP as a single row tensor.

\textbf{TinyMLP Inference Engine.}  The TinyMLP accepts the 5D
input and produces two parallel four-class logit vectors.
An $\operatorname{argmax}$ operation converts these to discrete
indices that select the next beacon interval from
$\mathcal{T} = \{0.1, 0.2, 0.5, 1.0\}$~s and the transmission
power offset from
$\mathcal{P} = \{-10, 0, +10, +20\}$~dBm, respectively.  The
entire inference executes in a single forward pass requiring no
iterative optimization at run time.

\textbf{ETSI Compliance Gate.}  The raw TinyMLP outputs are passed
through a hard clamp that enforces
$T_{\text{GenCam}} \in [100\,\text{ms},\, 1000\,\text{ms}]$ before
the timer is re-armed.  Since all four values in $\mathcal{T}$ already
lie within this range, the clamp acts as a safety guard against
implementation faults or adversarial inputs rather than a routine
correction step.  As noted above, any ETSI event trigger issued during
the interval overrides the timer and causes an immediate CAM
regardless of the scheduled $T_{\text{GenCam}}^{\mathrm{AI}}$ value.

\textbf{Learning Pipeline.}  The TinyMLP is trained entirely offline
through Behavior Cloning.  An \emph{oracle} is first constructed by
exhaustive grid-search over the $|\mathcal{T}| \times |\mathcal{P}| = 16$
action combinations for each sampled vehicle state; the oracle
selects the action that minimizes a weighted joint cost combining
AoI and CBR (detailed in Section~\ref{sec.prop}).  The resulting
state--action pairs form the Behavior Cloning dataset.  The TinyMLP
is trained on this dataset with a cross-entropy loss, converted to
ONNX format, post-training quantized to INT8, and exported as a
C header file (\texttt{tinymlp\_aidcc.h}) for on-board deployment.

\textbf{Evaluation Metrics.}  Performance is assessed using five
complementary metrics: (M1)~average Age of Information
($\overline{\text{AoI}}$, in ms), measuring information freshness
across all neighbor pairs; (M2)~Channel Busy Ratio (CBR), measuring
medium utilization; (M3)~Packet Delivery Ratio (PDR), measuring
reception success within 100~m; (M4)~energy efficiency (mJ/km),
measuring distance-normalized transmission energy; and
(M5)~ETSI compliance rate, measuring the fraction of CAM intervals
that respect the standard bounds.  Detailed definitions and
measurement procedures are given in Section~\ref{sec.s}.


% ============================================================
\section{Proposed Scheme: DQL-DCC}
\label{sec.prop}

Traditional heuristic models such as AdaptDCC exhibit inherent vulnerability
to local minima, oscillating between channel starvation and destructive bursts.
To overcome this, we propose DQL-DCC, an intelligent Deep Q-Learning (DQN)
based congestion controller integrated directly into the ETSI EN 302 637-2 MAC layer.
By casting the beacon scheduling process as a Markov Decision Process (MDP),
our agent autonomously learns a Pareto-optimal policy that maximizes Packet
Delivery Ratio (PDR) and minimizes energy consumption, while maintaining strict
compliance with ETSI CBR thresholds.

\subsection{System Architecture}
The DQL-DCC architecture replaces the conventional fixed-logic DCC state machine
with a deep neural network that evaluates the Q-values of available transmission
actions given the current vehicular context. The context is processed as a
three-dimensional state vector $\mathcal{S}$, which encapsulates the local
channel utilization, vehicle dynamics, and neighbor density. The network outputs
a joint action $A \in \mathcal{A}$ determining both the transmission interval
$T_{\text{GenCam}}$ and the transmit power $p_{\text{tx}}$.

\subsection{Algorithm Selection: Why DQN and MLP?}
\label{sec:justification}

A critical design choice in our architecture is the selection of the Deep Q-Network (DQN) algorithm coupled with a Multi-Layer Perceptron (MLP) over other modern Deep Learning or DRL alternatives.

\textbf{Why Reinforcement Learning?} Unlike supervised learning models (e.g., TabNet or ResNet) that require a ground-truth oracle, vehicular congestion control is a multi-agent NP-hard problem where a mathematical ``perfect action'' is unknown. DRL circumvents this by autonomously exploring the environment and discovering a Pareto-optimal policy maximizing delayed rewards.

\textbf{Why DQN over PPO, SAC, or TD3?} Advanced DRL algorithms such as Soft Actor-Critic (SAC) and Twin Delayed DDPG (TD3) are fundamentally designed for \emph{continuous} action spaces. However, ETSI EN 302 637-2 dictates a strictly \emph{discrete} set of permissible parameters for CAM generation. While Proximal Policy Optimization (PPO) supports discrete actions, it is an on-policy algorithm requiring vast amounts of fresh environmental interactions, rendering it highly sample-inefficient for computationally expensive SUMO simulations. DQN is natively discrete and off-policy, allowing highly efficient training via an experience replay buffer. Furthermore, Actor-Critic methods require complex network architectures, whereas DQN relies on a single lightweight Q-network during inference.

\textbf{Why an MLP over GRU or ResNet?} While sequence models like Gated Recurrent Units (GRU) excel at temporal dependencies, they require maintaining a continuously updated hidden state vector per vehicle, which severely burdens the limited memory footprint of a vehicular On-Board Unit (OBU). We bypass the need for RNNs by encapsulating temporal dynamics into a 1-second rolling average state variable ($\text{CBR}_{\text{smooth}}$). Consequently, a compact two-hidden-layer MLP (requiring under 4\,KB of MCU SRAM) is sufficient to process the 3-dimensional scalar state vector. This guarantees that the DQL-DCC inference completes in under 1\,ms on a standard ARM Cortex-M4 processor, strictly satisfying ETSI real-time scheduling constraints.


\subsection{Algorithm Selection: Why DQN and MLP?}
\label{sec:justification}

A critical design choice in our architecture is the selection of the Deep Q-Network (DQN) algorithm coupled with a Multi-Layer Perceptron (MLP) over other modern Deep Learning or DRL alternatives.

\textbf{Why Reinforcement Learning?} Unlike supervised learning models (e.g., TabNet or ResNet) that require a ground-truth oracle, vehicular congestion control is a multi-agent NP-hard problem where a mathematical ``perfect action'' is unknown. DRL circumvents this by autonomously exploring the environment and discovering a Pareto-optimal policy maximizing delayed rewards.

\textbf{Why DQN over PPO, SAC, or TD3?} Advanced DRL algorithms such as Soft Actor-Critic (SAC) and Twin Delayed DDPG (TD3) are fundamentally designed for \emph{continuous} action spaces. However, ETSI EN 302 637-2 dictates a strictly \emph{discrete} set of permissible parameters for CAM generation. While Proximal Policy Optimization (PPO) supports discrete actions, it is an on-policy algorithm requiring vast amounts of fresh environmental interactions, rendering it highly sample-inefficient for computationally expensive SUMO simulations. DQN is natively discrete and off-policy, allowing highly efficient training via an experience replay buffer. Furthermore, Actor-Critic methods require complex network architectures, whereas DQN relies on a single lightweight Q-network during inference.

\textbf{Why an MLP over GRU or ResNet?} While sequence models like Gated Recurrent Units (GRU) excel at temporal dependencies, they require maintaining a continuously updated hidden state vector per vehicle, which severely burdens the limited memory footprint of a vehicular On-Board Unit (OBU). We bypass the need for RNNs by encapsulating temporal dynamics into a 1-second rolling average state variable ($\text{CBR}_{\text{smooth}}$). Consequently, a compact two-hidden-layer MLP (requiring under 4\,KB of MCU SRAM) is sufficient to process the 3-dimensional scalar state vector. This guarantees that the DQL-DCC inference completes in under 1\,ms on a standard ARM Cortex-M4 processor, strictly satisfying ETSI real-time scheduling constraints.


\subsection{Markov Decision Process Formulation}
We formulate the vehicular transmission control problem as an MDP defined
by the tuple $(\mathcal{S}, \mathcal{A}, P, \mathcal{R}, \gamma)$.

\subsubsection{State Space $\mathcal{S}$}
The state space continuously tracks critical local network parameters. At any
decision step $t$, the state $s_t \in \mathcal{S}$ is defined as:
\begin{equation}
s_t = [ \text{CBR}_{\text{smooth}}, \| v \|, N_{\text{est}} ]
\end{equation}
where $\text{CBR}_{\text{smooth}}$ is the local 1-second rolling average of the
Channel Busy Ratio, $\| v \|$ is the vehicle's speed, and $N_{\text{est}}$ is the
estimated number of 1-hop neighbors derived from the most recent beacon history.

\subsubsection{Action Space $\mathcal{A}$}
The action space $\mathcal{A}$ defines the legal limits of CAM generation. 
To ensure strict adherence to ETSI specifications, the action space provides
discrete choices for the generation interval and transmission power:
\begin{equation}
a_t = [T_{\text{GenCam}}, p_{\text{tx}}]
\end{equation}
where $T_{\text{GenCam}} \in \{0.1, 0.2, 0.5, 1.0\}$ seconds, and 
$p_{\text{tx}} = 20.0$ dBm for maximum coverage. The agent selects an index 
$a_t \in \{0, 1, 2, 3\}$ mapping to the corresponding $T_{\text{GenCam}}$ value.

\subsubsection{Reward Function $\mathcal{R}$}
The reward function is the core driver of DQL-DCC's ability to maintain
unprecedented channel stability. To prevent the oscillatory bursts observed
in heuristic methods, we construct a heavily skewed reward function that
imposes a catastrophic penalty for CBR violations:
\begin{equation}
r_t = 
\begin{cases}
-100.0 \times (\text{CBR}_t - 0.60) & \text{if } \text{CBR}_t > 0.60 \\
0.1 \times N_{\text{sent}} & \text{otherwise}
\end{cases}
\end{equation}
This formulation forces the DRL agent to respect the ETSI 60\% CBR threshold
absolutely. The agent discovers that aggressively transmitting (low $T_{\text{GenCam}}$)
when the channel is near capacity yields catastrophic returns. Consequently,
it naturally converges to a conservative, highly stable transmission pattern that
virtually eliminates MAC layer collisions, maximizing the overall PDR.

\subsection{Deep Q-Network Training}
The DQL agent employs a Multi-Layer Perceptron (MLP) with two hidden layers
of 32 neurons each, employing ReLU activations. It is trained offline using
an $\epsilon$-greedy exploration strategy in a high-fidelity SUMO simulation
environment. The agent optimizes its policy by minimizing the Temporal Difference (TD)
loss using a target network and experience replay buffer. 
By learning from millions of simulated vehicular interactions, the network extracts
a generalized policy capable of scaling seamlessly from low-density highways to
hyper-congested urban intersections.


\section{Performance Evaluation}
\label{sec:evaluation}

The proposed DQL-DCC scheme is extensively evaluated and compared against
standard heuristic congestion control algorithms, including the ETSI
Reactive (ReactDCC) and Adaptive (AdaptDCC) approaches. We implement the
evaluation within a Python-based discrete-event simulation engine coupled
with the SUMO (Simulation of Urban MObility) traffic simulator.

\subsection{Channel Stability and Oscillation Prevention}
The most prominent failure mode of traditional heuristics like AdaptDCC is
their reliance on step-functions that cause massive oscillations in channel
utilization. When the channel is perceived as under-utilized, these heuristics
burst packets at maximum frequency, causing instantaneous CBR violations and
subsequent MAC-layer collisions.
In contrast, our DQL-DCC agent was strictly trained to penalize CBR violations.
As a result, DQL-DCC completely eradicates these channel bursts. Our time-series
analysis of instantaneous CBR demonstrates that while AdaptDCC oscillates wildly,
the DRL agent maintains a perfectly smooth and stable channel footprint, strictly
adhering to the 60\% ETSI threshold at all times.

\subsection{Packet Delivery Ratio (PDR) Maximization}
The primary consequence of eradicating channel bursts is a dramatic reduction in
MAC-layer collisions, which directly translates to unprecedented Packet Delivery Ratio
(PDR) performance.
Across all vehicle densities evaluated (from 10 to 100 vehicles per $km^2$), DQL-DCC
consistently achieves the highest PDR among all tested methods. For instance, in a
moderately dense scenario ($N=30$), DQL-DCC achieves a PDR of 58.6\%, significantly
outperforming AdaptDCC (54.0\%) and ReactDCC (53.7\%). This confirms that the DRL
agent successfully learned a Pareto-optimal transmission scheduling policy that avoids
destructive interference.

\subsection{Energy Efficiency}
In addition to maximizing reliability, DQL-DCC achieves state-of-the-art energy
efficiency. By refraining from unnecessary, collision-prone high-frequency bursting,
the agent drastically reduces the total number of CAMs transmitted per vehicle.
At $N=30$, DQL-DCC consumes only 4.53~mJ on average, marking a near 20\% energy saving
compared to AdaptDCC (5.66~mJ). This makes DQL-DCC highly suitable for resource-constrained
vehicular IoT modules where battery and thermal budgets are strict constraints.

\subsection{The AoI Trade-off}
The strict standard-compliant stability and high reliability of DQL-DCC comes at a slight,
yet theoretically necessary, trade-off in Age of Information (AoI). Heuristic methods
that momentarily violate CBR limits can achieve superficially lower average AoI values
due to the illegal burst transmissions. However, as demonstrated by the PDR results,
these lower AoI packets are frequently lost to collisions. DQL-DCC sacrifices a fraction
of AoI freshness (operating in the 1000ms--1400ms range at extreme densities) to guarantee
that when a packet is sent, it is reliably delivered. In safety-critical V2X applications,
guaranteed delivery (PDR) under strict standard compliance is vastly superior to sporadic,
collision-heavy bursts of fresh information.


\section{Conclusion}
\label{sec.c}
% ============================================================

This paper addressed the fundamental tension between information
freshness and channel congestion in ETSI CAM-based vehicular IoT
networks.  Standard DCC algorithms, being reactive and CBR-centric,
cannot jointly optimize the Age of Information and channel utilization,
nor do they leverage the kinematic context of individual vehicles.
We proposed \textbf{TinyMLP-AI-DCC}, a two-layer MLP with 1,448
parameters that accepts a five-dimensional vehicle state vector
(speed, acceleration, estimated neighbor count, CBR, and AoI) and
outputs a joint action $(T_{\text{GenCam}}, p_{\text{tx}})$ in a
single forward pass.  The model is trained offline via Behavior
Cloning against a myopically optimal oracle generated by exhaustive
16-point grid search, avoiding the convergence instability of deep
reinforcement learning.  Post-training INT8 quantization reduces the
deployment footprint to under 4~KB of MCU memory---fitting
comfortably within the 192~KB SRAM of an STM32F407 (ARM Cortex-M4)
and completing inference in under 1~ms at 80~MHz.

The baseline comparison study reported in this work provides
quantitative evidence for the AoI--CBR Pareto gap that motivates the
proposed design.  Across the \texttt{urban\_grid} scenario (three
random seeds, $\sim$88 concurrent vehicles, density $\approx$
3.06~veh/(km$\cdot$lane)), the four baseline methods reveal a clear
frontier: Fixed~10~Hz (BL-D) achieves the lowest AoI of 321.96~ms
but pushes CBR to 0.6208 (near-overload); Variable Beacon (BL-C)
achieves the lowest CBR of 0.4098 but at an AoI cost of 501.41~ms;
while the two ETSI-standard variants (BL-A/B) occupy a moderate but
unoptimized middle ground.  No existing baseline simultaneously
achieves low AoI \emph{and} moderate CBR, confirming the existence of
an open Pareto-dominant operating region that the TinyMLP-AI-DCC
controller is designed to target, in direct correspondence with the
joint cost $J(T, p)$ defined in Section~\ref{sec.prop}.

\textbf{Limitations.}  The present version of this paper has two
principal limitations that must be acknowledged.  First, the
end-to-end evaluation of the proposed TinyMLP-AI-DCC controller---its
training, BC policy convergence, and five-metric comparison against
the four baselines---has not yet been completed, and all entries in
Table~\ref{tab:baseline_main} for the Proposed method remain pending.
Consequently, the performance claims in the abstract (``up to 35\%
AoI reduction'') are projections based on the oracle performance
analysis rather than empirical end-to-end results.  Second, the MCU
resource figures (inference latency and memory footprint) are
simulation estimates derived from publicly available CMSIS-NN
benchmarks; no physical measurement on actual STM32F407 hardware has
been performed in this version.

\textbf{Future Work.}
The immediate next steps are: (i)~completion of the TinyMLP Behavior
Cloning training pipeline and end-to-end simulation evaluation;
(ii)~ablation studies (ABL-1: rate-only, ABL-2: no-AoI input) and
sensitivity sweeps (SA1--SA4) as detailed in
Section~\ref{sec.s_pending}; (iii)~extension to the highway mobility
scenario~(S2) and a density sweep across 10--30~veh/(km$\cdot$lane)
to characterise the controller's generalization envelope; and
(iv)~physical deployment on an actual STM32F407 board to obtain
measured latency and memory consumption data, replacing the current
simulation-based estimates and providing a concrete, reproducible MCU
benchmark.

% ============================================================
\begin{thebibliography}{00}
% ********** Related Work: AI-driven MAC Protocols (Direct competition) **********
% (10 papers)
\bibitem{Wu2025} Chien-Min Wu, Cheng-Tai Tsai, Cheng-Chun Hou, \emph{et al.}, ``Emergency Message Broadcast Mechanism in Vehicular Ad-Hoc Networks Based on Reinforcement Learning With Contention Estimation,'' \emph{IEEE Transactions on Intelligent Vehicles}, 2025. doi: 10.1109/TIV.2024.3418778.
\bibitem{Iqbal2025} Adeel Iqbal, A. Nauman, Tahir Khurshaid, ``Lightweight Reinforcement Learning for Priority-Aware Spectrum Management in Vehicular IoT Networks,'' \emph{MDPI Sensors}, 2025. doi: 10.3390/s25216777.
\bibitem{Louvros2026} Spyridon Louvros, Anupkumar Pandey, B. Shah, Yashesh Buch, ``Federated Edge Learning MAC Cross-Layer Optimization in 6G Networks : The Case for Radial Basis Functions Optimization,'' \emph{2026 Panhellenic Conference on Electronics & Telecommunications (PACET)}, 2026. doi: 10.1109/PACET68758.2026.11498274.
\bibitem{Li2025a} Zihan Li, Ping Wang, Yamin Shen, Song Li, ``Reinforcement Learning-Based Resource Allocation Scheme of NR-V2X Sidelink for Joint Communication and Sensing,'' \emph{MDPI Sensors}, 2025. doi: 10.3390/s25020302.
\bibitem{Wang2024} Ping Wang, C. Zheng, Chenggen Pu, Zhiyi Tian, ``Dynamic Multi-Channel Access in Heterogeneous Industrial Wireless Sensor Networks: A MADRL Approach,'' \emph{ACM Cloud and Autonomic Computing Conference}, 2024. doi: 10.1109/CAC63892.2024.10865069.
\bibitem{Ibrahim2025} Amin S. Ibrahim, Ahmed M. Abbas, Saeed Mohsen, ``A novel adaptive IEEE 802.11ah MAC protocol using machine learning for data collision reduction in smart cities IoT networks,'' \emph{Wireless networks}, 2025. doi: 10.1007/s11276-025-04059-2.
\bibitem{Ibrahim2026} Amin S. Ibrahim, Ahmed M. Abbas, Saeedi Mohsen, ``Correction: A novel adaptive IEEE 802.11ah MAC protocol using machine learning for data collision reduction in smart cities IoT networks,'' \emph{Wireless networks}, 2026. doi: 10.1007/s11276-025-04078-z.
\bibitem{Zila2026} Amine Zila, Y. Mouzouna, A. Ouchatti, Ikram Daanoune, ``Edge AI and TinyML for Enhancing MAC Protocols: A New Paradigm for Wireless Sensor Networks in IIoT,'' \emph{International Journal of Communication Systems}, 2026. doi: 10.1002/dac.70403.
\bibitem{Ni2024} Yang Ni, Danny Abraham, Mariam Issa, \emph{et al.}, ``Dynamic MAC Protocol for Wireless Spectrum Sharing via Hyperdimensional Self-Learning,'' \emph{IEEE Access}, 2024. doi: 10.1109/ACCESS.2024.3464868.
\bibitem{PV2025} I. P., T. Venkatesh, ``RIS-Assisted Full Duplex MAC Protocol for NR-V2X with ML-Driven QoS Prediction,'' \emph{2025 IEEE Future Networks World Forum (FNWF)}, 2025. doi: 10.1109/FNWF66845.2025.11317280.

% ********** Related Work: Lightweight ML for Vehicular Networks **********
% (13 papers)
\bibitem{Li2025} Junjun Li, Yanyan Ma, Jiahui Bai, \emph{et al.}, ``A Lightweight Intrusion Detection System with Dynamic Feature Fusion Federated Learning for Vehicular Network Security,'' \emph{MDPI Sensors}, 2025. doi: 10.3390/s25154622.
\bibitem{Yan2025} Na Yan, Senior Member Ieee Kezhi Wang, Kangda Zhi, \emph{et al.}, ``Secure and Private Over-the-Air Federated Learning: Biased and Unbiased Aggregation Design,'' \emph{IEEE Transactions on Wireless Communications}, 2025. doi: 10.1109/TWC.2025.3550159.
\bibitem{Liu2024} Yiming Liu, Qiang Wang, Jiaao Chen, Wenqi Zhang, Chen Sun, ``Federated Multi-Agent Deep Reinforcement Learning Approach for Resource Allocation in Platoon-Based NR-V2X,'' \emph{IEEE Vehicular Technology Conference}, 2024. doi: 10.1109/VTC2024-Spring62846.2024.10683498.
\bibitem{Wu2024} Maoqiang Wu, Guoliang Cheng, Dongdong Ye, \emph{et al.}, ``Federated Split Learning With Data and Label Privacy Preservation in Vehicular Networks,'' \emph{IEEE Transactions on Vehicular Technology}, 2024. doi: 10.1109/TVT.2023.3304176.
\bibitem{Zhang2025} Jianquan Zhang, Fangting Huang, Shuqing Zhu, Xiao Xiao, ``A Resource Allocation Strategy in Internet of Vehicles Based on Multi-Task Federated Learning and Incentive Mechanism,'' \emph{IEEE transactions on intelligent transportation systems (Print)}, 2025. doi: 10.1109/TITS.2025.3528969.
\bibitem{Yang2024a} Lu Yang, Songtao Guo, Chen-Khong Tham, \emph{et al.}, ``PreM-FedIoV: A Novel Federated Reinforcement Learning Framework for Predictive Maintenance in IoV,'' \emph{IEEE Transactions on Mobile Computing}, 2024. doi: 10.1109/TMC.2024.3404249.
\bibitem{Xu2024} Mengfan Xu, Yaguang Lin, ``FedSteg: Coverless Steganography‐Based Privacy‐Preserving Decentralized Federated Learning,'' \emph{IEEJ TRANSACTIONS ON ELECTRICAL AND ELECTRONIC ENGINEERING}, 2024. doi: 10.1002/tee.24085.
\bibitem{Elzemity2024} Adel Elzemity, Budi Arief, ``Privacy Threats and Countermeasures in Federated Learning for Internet of Things: A Systematic Review,'' \emph{2024 IEEE International Conferences on Internet of Things (iThings) and IEEE Green Computing & Communications (GreenCom) and IEEE Cyber, Physical & Social Computing (CPSCom) and IEEE Smart Data (SmartData) and IEEE Congress on Cybermatics}, 2024. doi: 10.1109/iThings-GreenCom-CPSCom-SmartData-Cybermatics62450.2024.00072.
\bibitem{Liu2025} Xianhui Liu, Jianle Liu, Yingyao Zhang, Ning Jia, Chenlin Zhu, ``Heterogeneous Federated Learning for Vehicle-to-Everything: Feature Prototype Aggregation and Generative Feedback Mechanism,'' \emph{IEEE Open Journal of Vehicular Technology}, 2025. doi: 10.1109/OJVT.2025.3594030.
\bibitem{Zhang2025b} Zhiwei Zhang, ``A Vehicular Network Intrusion Detection Method Based on Knowledge Distillation and Hybrid Learning,'' \emph{2025 International Conference on Algorithms, Data Mining, and Information Technology (ADMIT)}, 2025. doi: 10.1109/ADMIT67050.2025.11337192.
\bibitem{Liao2025} Gengjian Liao, Yulin Yang, Zhenni Feng, ``Personalized federated learning through self-knowledge distillation in vehicular edge computing,'' \emph{Comput. Networks}, 2025. doi: 10.1016/j.comnet.2025.111406.
\bibitem{Li2025b} Chunlin Li, Yong Zhang, Long Yu, Mengjie Yang, ``Efficient Vehicle Selection and Resource Allocation for Knowledge Distillation-Based Federated Learning in UAV-Assisted VEC,'' \emph{IEEE transactions on intelligent transportation systems (Print)}, 2025. doi: 10.1109/TITS.2025.3525735.
\bibitem{Zhao2026} Pengcheng Zhao, Xingjia Wei, Yong Wang, Xin Li, Shibao Sun, ``AFed-CHKD: An Adversarial Federated Learning-Based Cluster Head Selection and Knowledge Distillation in Vehicular Sensor Network,'' \emph{IEEE Sensors Journal}, 2026. doi: 10.1109/JSEN.2026.3665994.

% ********** Related Work: V2X / NR-V2X Protocols **********
% (10 papers)
\bibitem{Nikooroo2024} Saleh Nikooroo, Juan Estrada-Jimenez, Aurel Machalek, \emph{et al.}, ``Poster: Evaluating the Potential of Mode 2(b) Resource Allocation in NR V2X Sidelink,'' \emph{IEEE Vehicular Networking Conference}, 2024. doi: 10.1109/VNC61989.2024.10575999.
\bibitem{Li2024} Pengfei Li, Xin-Lin Huang, Fei Hu, ``When DSRC Meets With C-V2X: How to Jointly Predict and Select Idle Channels?,'' \emph{IEEE Transactions on Vehicular Technology}, 2024. doi: 10.1109/TVT.2024.3399211.
\bibitem{Gao2024} Ang Gao, Ziqing Zhu, Jiankang Zhang, Wei Liang, Yansu Hu, ``Matching Combined Heterogeneous Multi-Agent Reinforcement Learning for Resource Allocation in NOMA-V2X Networks,'' \emph{IEEE Transactions on Vehicular Technology}, 2024. doi: 10.1109/TVT.2024.3409048.
\bibitem{Narayanasamy2024} Iswarya Narayanasamy, Venkateswari Rajamanickam, ``A Cascaded Multi-Agent Reinforcement Learning-Based Resource Allocation for Cellular-V2X Vehicular Platooning Networks,'' \emph{MDPI Sensors}, 2024. doi: 10.3390/s24175658.
\bibitem{Chen2025} Chen Chen, Wenhua Wang, Ziye Liu, \emph{et al.}, ``RLFN-VRA: Reinforcement Learning-Based Flexible Numerology V2V Resource Allocation for 5G NR V2X Networks,'' \emph{IEEE Transactions on Intelligent Vehicles}, 2025. doi: 10.1109/TIV.2024.3427399.
\bibitem{Mao2024} Yaqi Mao, Xin Yang, Ling Wang, \emph{et al.}, ``A High-Capacity MAC Protocol for UAV-Enhanced RIS-Assisted V2X Architecture in 3-D IoT Traffic,'' \emph{IEEE Internet of Things Journal}, 2024. doi: 10.1109/JIOT.2024.3387997.
\bibitem{Iliopoulos2025} Christos Iliopoulos, Athanasios C. Iossifides, C. Foh, Periklis Chatzimisios, ``IEEE 802.11BD for Next-Generation V2X Communications: From Protocol to Services,'' \emph{IEEE Communications Standards Magazine}, 2025. doi: 10.1109/MCOMSTD.2025.3569015.
\bibitem{Zhang2025a} Zhishuo Zhang, Wen Huang, Ying Huang, \emph{et al.}, ``An Attribute-Based Pre-Authenticated Secure Communication Protocol Enabling Key Protection and Credential Online-Upgrading for 5G NR V2X,'' \emph{IEEE transactions on intelligent transportation systems (Print)}, 2025. doi: 10.1109/TITS.2025.3558978.
\bibitem{Alagumani2024} Sangeetha Alagumani, Umamaheswari Natarajan, ``Q-learning and fuzzy logic multi-tier multi-access edge clustering for 5g v2x communication,'' \emph{Network}, 2024. doi: 10.1080/0954898X.2024.2309947.
\bibitem{Doğan2026} Yasemin Doğan, Ramazan Ünlü, ``Machine Learning for V2X-Enabled Microgrids: A Bibliometric and Thematic Review of Intelligent Energy Management Applications,'' \emph{The Arabian journal for science and engineering}, 2026. doi: 10.1007/s13369-026-11201-5.

% ********** Background: Vehicular Networking, Edge Intelligence **********
% (14 papers)
\bibitem{Rawlley2026} Oshin Rawlley, Shashank Gupta, Jatin Kumar Panwar, Palak Sharma, Shailendra Rathore, ``Asynchronous Deep Reinforcement Learning for Semantic Communication and Digital-Twin Deployment in Transportation Networks,'' \emph{IEEE transactions on intelligent transportation systems (Print)}, 2026. doi: 10.1109/TITS.2025.3593038.
\bibitem{Mianji2025} Elham Mohammadzadeh Mianji, Gabriel-Miro Muntean, Irina Tal, ``Enhancing Vehicular Network Security, Privacy, and Trust Through Reinforcement Learning: A Comprehensive Survey,'' \emph{IEEE transactions on intelligent transportation systems (Print)}, 2025. doi: 10.1109/TITS.2025.3612202.
\bibitem{Yang2024} Liu Yang, Yifei Wei, Zhiyong Feng, Qixun Zhang, Zhu Han, ``Deep Reinforcement Learning-Based Resource Allocation for Integrated Sensing, Communication, and Computation in Vehicular Network,'' \emph{IEEE Transactions on Wireless Communications}, 2024. doi: 10.1109/TWC.2024.3470873.
\bibitem{Geraci2025} Giovanni Geraci, Francesca Meneghello, F. Wilhelmi, \emph{et al.}, ``Wi-Fi: 25 Years and Counting,'' \emph{Proceedings of the IEEE}, 2025. doi: 10.1109/JPROC.2026.3662513.
\bibitem{Xie2025} Xin Xie, Tao Zhong, Heng Wang, ``Scheduling for Maximizing the Information Freshness in Vehicular Edge Computing- Assisted IoT Systems,'' \emph{IEEE transactions on intelligent transportation systems (Print)}, 2025. doi: 10.1109/TITS.2024.3514099.
\bibitem{Chen2026} Xiangyi Chen, Huanlai Xing, Qinnan Zhang, \emph{et al.}, ``Age-of-Information-Aware Hierarchical Collaborative Inference in Digital Twin-Enabled Vehicular Edge Computing,'' \emph{IEEE Transactions on Network Science and Engineering}, 2026. doi: 10.1109/TNSE.2025.3615533.
\bibitem{Liu2024a} Jianhang Liu, Kunlei Xue, Q. Miao, \emph{et al.}, ``MCVCO: Multi-MEC Cooperative Vehicular Computation Offloading,'' \emph{IEEE Transactions on Intelligent Vehicles}, 2024. doi: 10.1109/TIV.2023.3299381.
\bibitem{Huang2025} Mengting Huang, Wenqian Zhang, Guanglin Zhang, ``Joint Service Placement and Task Offloading in Vehicle-Edge-Cloud Collaborative Networks,'' \emph{IEEE transactions on intelligent transportation systems (Print)}, 2025. doi: 10.1109/TITS.2025.3590099.
\bibitem{Zhou2025} Huan Zhou, Wenchao Xia, Gan Zheng, Hongbo Zhu, ``Graph Neural Network-Based Continual Learning for Resource Allocation in Dynamic Wireless Environments,'' \emph{IEEE Transactions on Vehicular Technology}, 2025. doi: 10.1109/TVT.2025.3585146.
\bibitem{Li2024a} Mengfei Li, Wei Wang, Shijie Feng, \emph{et al.}, ``GCN-Enhanced Multi-Agent MAC Protocol for Vehicular Communications,'' \emph{International Conference on Industrial Informatics}, 2024. doi: 10.1109/INDIN58382.2024.10774514.
\bibitem{Tang2024} Zhiqing Tang, Fangyi Mou, Jiong Lou, \emph{et al.}, ``Multi-User Layer-Aware Online Container Migration in Edge-Assisted Vehicular Networks,'' \emph{IEEE/ACM Transactions on Networking}, 2024. doi: 10.1109/TNET.2023.3330255.
\bibitem{Bhattacharyya2024} Sagnik Bhattacharyya, Pankaj Kumar, Sam Darshi, \emph{et al.}, ``Hybrid Relaying Based Cross Layer MAC Protocol Using Variable Beacon for Cooperative Vehicles,'' \emph{IEEE Transactions on Vehicular Technology}, 2024. doi: 10.1109/TVT.2023.3307672.
\bibitem{Atimati2024} Ehinomen Atimati, David H. Crawford, R. Stewart, ``Intelligent Joint Resource Management of Shared Spectrum HetNets,'' \emph{2024 IEEE Future Networks World Forum (FNWF)}, 2024. doi: 10.1109/FNWF63303.2024.11028798.
\bibitem{Alsadie2024} Deafallah Alsadie, ``Artificial Intelligence Techniques for Securing Fog Computing Environments: Trends, Challenges, and Future Directions,'' \emph{IEEE Access}, 2024. doi: 10.1109/ACCESS.2024.3463791.

\end{thebibliography}
% ============================================================

\end{document}
